% GAL1897_W07 — English translation of physical PDF pages 044–053.
% Primary authority: the frozen 1897 French diplomatic transcription from
% JP2 leaves L0043–L0052, navigated against the canonical 96-page 1897 PDF.
% Definite source-level mathematical defects are translated without silent repair
% and identified in adjacent critical notes from the frozen critical baseline.
% PDF 053 / printed page 24 remains the inspected blank-topology leaf.

\providecommand{\GalSourcePageStart}[4]{}%
\providecommand{\GalSourcePageBoundary}[4]{\space}%
\providecommand{\GalSourceBlankPage}[4]{}%
\providecommand{\GalSourceError}[2]{#1}%
\providecommand{\GalPrintedError}[2]{#1}%
\providecommand{\GalHistoricalFootnoteText}[2]{\footnotetext{#2}}%
\providecommand{\GalHistoricalFootnote}[2]{\footnote{#2}}%
\providecommand{\GalSourceImage}[5]{}%
\providecommand{\GalArticleEndOrnament}[1]{\par\smallskip\begin{center}\rule{0.18\linewidth}{0.4pt}\end{center}}%

\GalSourcePageStart{044}{L0043}{15}{15}

% ENSEG:W07:PDF044:0001
\begin{center}
{\scriptsize\bfseries ON\par}
\vspace{0.25em}
{\large\bfseries THE THEORY OF NUMBERS\textsuperscript{(*)}.\par}
\vspace{0.55em}
\rule{0.12\linewidth}{0.4pt}
\end{center}

% ENSEG:W07:PDF044:0002
\GalHistoricalFootnoteText{*}{M.~Férussac's \emph{Bulletin des Sciences mathématiques}, vol.~XIII,
p.~438 (1830, June issue), with the following note: ``This Memoir forms part
of M.~Galois's research on the theory of permutations and algebraic
equations.'' \textsc{(J. Liouville.)}}

% ENSEG:W07:PDF044:0003
When one agrees to regard as zero all quantities that, in algebraic
calculations, are found multiplied by a given prime number $p$, and,
under this convention, seeks the solutions of an algebraic equation
$Fx=0$, which M.~Gauss denotes by $Fx\equiv0$, it is customary to
consider only the integral solutions of such questions. Having been led
by particular investigations to consider incommensurable solutions, I
have obtained some results that I believe to be new.

% ENSEG:W07:PDF044:0004
Let there be such an equation or congruence, $Fx=0$, with modulus $p$.
Suppose first, for greater simplicity, that the congruence in question
admits no commensurable factor; that is, that one cannot find three
functions $\varphi x$, $\psi x$, and $\chi x$ such that
% ENSEG:W07:PDF044:0005
\[
\varphi x\mathbin{\cdot}\psi x=Fx+p\mathbin{\cdot}\chi x.
\]
% ENSEG:W07:PDF044:0006
In this case the congruence will therefore admit no integral root, nor
even any incommensurable root of lower degree. The roots of this
congruence must therefore be regarded as a kind of imaginary symbol,
since they do not satisfy questions concerning the integers; the use of
these symbols in calculation will often be as useful as that of the
imaginary $\sqrt{-1}$ in ordinary analysis.

% ENSEG:W07:PDF044:0007
We shall be concerned with classifying these imaginaries and reducing
them to the smallest possible number.

% ENSEG:W07:PDF044:0008
Call $i$ one of the roots of the congruence $Fx=0$, which we suppose
to have degree $\nu$.

% ENSEG:W07:PDF044:0009
Consider the general expression
% ENSEG:W07:PDF044:0010
\[
\tag{A}
a+a_1i+a_2i^2+\cdots+a_{\nu-1}i^{\nu-1},
\]

\GalSourcePageBoundary{045}{L0044}{16}{16}

% ENSEG:W07:PDF045:0011
where $a$, $a_1$, $a_2$, \ldots, $a_{\nu-1}$ represent integers. When
these numbers are given all possible values, expression (A) takes
$p^\nu$ values, which, as I shall show, possess the same properties as
the natural numbers in the \emph{theory of power residues}.

% ENSEG:W07:PDF045:0012
Of the expressions (A), take only the $p^\nu-1$ values for which $a$,
$a_1$, $a_2$, \ldots, $a_{\nu-1}$ are not all zero; let $\alpha$ be one
of these expressions.

% ENSEG:W07:PDF045:0013
If $\alpha$ is raised successively to the second, third, \ldots,
powers, one obtains a sequence of quantities of the same form [because
every function of $i$ can be reduced to degree $\nu-1$]. Thus one must
have $\alpha^n=1$ for some number $n$; let $n$ be the least number for
which $\alpha^n=1$. There will be a set of $n$ mutually distinct
expressions,
% ENSEG:W07:PDF045:0014
\[
1,\ \alpha,\ \alpha^2,\ \alpha^3,\ \ldots,\ \alpha^{n-1}.
\]
% ENSEG:W07:PDF045:0015
Multiply these $n$ quantities by another expression $\beta$ of the same
form. We obtain another new group of quantities, all different from the
first and from one another. If the quantities (A) have not been
exhausted, multiply the powers of $\alpha$ again by a new expression
$\gamma$, and so on. One thus sees that the number $n$ must divide the
total number of quantities (A). Since this number is $p^\nu-1$, one sees
that $n$ divides $p^\nu-1$. It follows further that
% ENSEG:W07:PDF045:0016
\[
\alpha^{p^\nu-1}=1,\qquad\text{or}\qquad \alpha^{p^\nu}=\alpha.
\]

% ENSEG:W07:PDF045:0017
Next one proves, as in the theory of numbers, that there are primitive
roots $\alpha$ for which precisely $p^\nu-1=n$, and which consequently
reproduce, on being raised to powers, the entire sequence of the other
roots.

% ENSEG:W07:PDF045:0018
And any one of these primitive roots will depend only on a congruence
of degree $\nu$, an \emph{irreducible} congruence; otherwise the equation
in $i$ would not be irreducible either, because the roots of the
congruence in $i$ are all powers of the primitive root.

% ENSEG:W07:PDF045:0019
We see here the remarkable consequence that all algebraic quantities
that can arise in the theory are roots of equations of the form
% ENSEG:W07:PDF045:0020
\[
x^{p^\nu}=x.
\]

\GalSourcePageBoundary{046}{L0045}{17}{17}

% ENSEG:W07:PDF046:0021
Stated algebraically, this proposition is as follows: given a function
$Fx$ and a prime number $p$, one can write
% ENSEG:W07:PDF046:0022
\[
fx\mathbin{\cdot}Fx=x^{p^\nu}-x+p\varphi x,
\]
% ENSEG:W07:PDF046:0023
where $fx$ and $\varphi x$ are integral functions, whenever the
congruence $Fx=0\pmod p$ is irreducible.

% ENSEG:W07:PDF046:0024
If one wishes to obtain all the roots of such a congruence from a single
one, it is enough to observe that, in general,
% ENSEG:W07:PDF046:0025
\[
(Fx)^{p^n}=F\bigl(x^{p^n}\bigr)
\]
% ENSEG:W07:PDF046:0026
and consequently, if one root is $x$, the others will be
% ENSEG:W07:PDF046:0027
\[
x^p,\ x^{p^2},\ \ldots,\ x^{p^{\nu-1}}
% ENSEG:W07:PDF046:0028
\GalHistoricalFootnote{1}{From the fact that the roots of the irreducible congruence of degree $\nu$
% ENSEG:W07:PDF046:0029
\[
Fx=0
\]
are expressed by the sequence
% ENSEG:W07:PDF046:0030
\[
x,\ x^p,\ x^{p^2},\ \ldots,\ x^{p^{\nu-1}},
\]
\GalSourceError{it would be wrong to conclude that these roots are always quantities
expressible by radicals. Here is an example to the contrary:}{W07-SE001: the displayed example does not establish non-radical expressibility; its roots are nontrivial cube roots of unity in F_4.}

The irreducible congruence
% ENSEG:W07:PDF046:0031
\[
x^2+x+1=0\pmod 2
\]
gives
% ENSEG:W07:PDF046:0032
\[
x=\frac{-1+\sqrt{-3}}{2},
\]
which reduces to
% ENSEG:W07:PDF046:0033
\[
\frac{0}{0}\pmod 2,
\]
a formula that tells us nothing.}.
\]

% ENSEG:W07:PDF046:0034
\GalCriticalNote{W07-SE001}{The quadratic example is still solvable by radicals}{The
historical footnote treats $x^2+x+1=0\pmod 2$ as an example whose roots cannot
be expressed by radicals because the ordinary quadratic formula becomes
$0/0$. But a root $\alpha$ satisfies $\alpha^3=1$ and is itself a nontrivial
cube root of unity in $\mathbf F_4$. The indeterminate quotient shows only
that the characteristic-zero formula has failed; it does not prove
non-radicality. The footnote is retained, but cannot serve as a counterexample.}

% ENSEG:W07:PDF046:0035
It remains to show that, conversely to what we have just said, all the
roots of the equation or congruence $x^{p^\nu}=x$ depend on a single
congruence of degree $\nu$.

% ENSEG:W07:PDF046:0036
Indeed, let $i$ be a root of an irreducible congruence such that all the
roots of the congruence $x^{p^\nu}=x$ are rational functions of $i$ (it
is clear that this property holds here as it does for ordinary
equations)
% ENSEG:W07:PDF046:0122
\GalHistoricalFootnote{2}{The general proposition in question can be stated as follows: given an
algebraic equation, one can find a rational function $\theta$ of all its
roots such that, conversely, each root is expressed rationally in terms of
$\theta$. This theorem was known to Abel, as may be seen from the first Part
of the Memoir left by that celebrated geometer on elliptic functions.

\hfill E. G.}.

\GalSourcePageBoundary{047}{L0046}{18}{18}

% ENSEG:W07:PDF047:0037
It is first evident that the degree $\mu$ of the congruence in $i$ cannot
be less than $\nu$, for otherwise the congruence
% ENSEG:W07:PDF047:0038
\[
\tag{$\nu$}
x^{p^\nu-1}-1=0
\]
% ENSEG:W07:PDF047:0039
would have all its roots in common with the congruence
% ENSEG:W07:PDF047:0040
\[
x^{p^\mu-1}-1=0,
\]
% ENSEG:W07:PDF047:0041
which is absurd, since congruence $(\nu)$ has no equal roots, as one sees
by taking the derivative of its left-hand side. I now say that $\mu$
cannot be greater than $\nu$ either.

% ENSEG:W07:PDF047:0042
Indeed, if that were so, all the roots of the congruence
% ENSEG:W07:PDF047:0043
\[
x^{p^\mu}=x
\]
% ENSEG:W07:PDF047:0044
would have to depend rationally on those of the congruence
% ENSEG:W07:PDF047:0045
\[
x^{p^\nu}=x.
\]
% ENSEG:W07:PDF047:0046
But it is easy to see that, if
% ENSEG:W07:PDF047:0047
\[
i^{p^\nu}=i,
\]
% ENSEG:W07:PDF047:0048
every rational function $h=fi$ will also give
% ENSEG:W07:PDF047:0049
\[
(fi)^{p^\nu}=f\bigl(i^{p^\nu}\bigr)=fi,\qquad\text{whence}\qquad h^{p^\nu}=h.
\]

% ENSEG:W07:PDF047:0050
Thus all the roots of the congruence $x^{p^\mu}=x$ would be common
to it and to the equation $x^{p^\nu}=x$, which is absurd.

% ENSEG:W07:PDF047:0124
\GalCriticalNote{POST-P13-A004}{Why the two degrees are equal}{The constructed
field $\mathbf F_{p^\nu}$ lies in $\mathbf F_{p^\mu}$, so $\nu\mid\mu$.
Conversely, because $i$ is rational in the constructed quantity,
$\mathbf F_{p^\mu}$ lies in $\mathbf F_{p^\nu}$, so $\mu\mid\nu$. Hence
$\mu=\nu$, and the compressed argument is sound.}

% ENSEG:W07:PDF047:0051
We now know, therefore, that all the roots of the equation or
congruence $x^{p^\nu}=x$ necessarily depend on a \emph{single
irreducible} congruence of degree $\nu$.

% ENSEG:W07:PDF047:0052
Now, to obtain this irreducible congruence on which the roots of the
congruence $x^{p^\nu}=x$ depend, the most general method will be first
to free this congruence of all factors it may have in common with
congruences of lower degree and of the form
% ENSEG:W07:PDF047:0053
\[
x^{p^\mu}=x.
\]

% ENSEG:W07:PDF047:0054
One will thereby obtain a congruence that must split into irreducible
congruences of degree $\nu$. And since one knows how to express all the
roots of each of these irreducible congruences in terms of a single one,

\GalSourcePageBoundary{048}{L0047}{19}{19}

% ENSEG:W07:PDF048:0055
it will be easy to obtain them all by M.~Gauss's method.

% ENSEG:W07:PDF048:0056
Most often, however, it will be easy by trial to find an irreducible
congruence of a given degree $\nu$, and all the others should be deduced
from it.

% ENSEG:W07:PDF048:0057
As an example, let $p=7$, $\nu=3$. Let us seek the roots of the
congruence
% ENSEG:W07:PDF048:0058
\[
\tag{1}
x^{7^3}=x\pmod 7.
\]
% ENSEG:W07:PDF048:0059
I observe that the congruence
% ENSEG:W07:PDF048:0060
\[
\tag{2}
i^3=2\pmod 7
\]
% ENSEG:W07:PDF048:0061
is irreducible and of degree $3$, and all the roots of congruence (1)
therefore depend rationally on those of congruence (2), so that all the
roots of (1) are of the form
% ENSEG:W07:PDF048:0062
\[
\tag{3}
a+a_1i+a_2i^2\qquad\text{or}\qquad
a+a_1\sqrt[3]{2}+a_2\sqrt[3]{4}.
\]

% ENSEG:W07:PDF048:0063
We must now find a primitive root; that is, a form of expression (3)
which, when raised to all powers, gives all the roots of the congruence
% ENSEG:W07:PDF048:0064
\[
x^{7^3-1}=1,\qquad\text{namely}\qquad
x^{2^1\mathbin{\cdot}3^2\mathbin{\cdot}19}=1\pmod 7,
\]
% ENSEG:W07:PDF048:0065
and for this we need only obtain a primitive root of each congruence
% ENSEG:W07:PDF048:0066
\[
x^2=1,\qquad x^{3^2}=1,\qquad x^{19}=1.
\]

% ENSEG:W07:PDF048:0067
The primitive root of the first is $-1$; those of $x^{3^2}-1=0$ are
given by the equations
% ENSEG:W07:PDF048:0068
\[
x^3=2,\qquad x^3=4,
\]
% ENSEG:W07:PDF048:0069
so that $i$ is a primitive root of $x^{3^2}=1$.

% ENSEG:W07:PDF048:0070
It remains only to find a root of $x^{19}-1=0$, or rather of
% ENSEG:W07:PDF048:0071
\[
\frac{x^{19}-1}{x-1}=0,
\]
% ENSEG:W07:PDF048:0072
and for this let us try whether the question can be satisfied

\GalSourcePageBoundary{049}{L0048}{20}{20}

% ENSEG:W07:PDF049:0073
by simply setting $x=a+a_1i$ instead of $a+a_1i+a_2i^2$. We must have
% ENSEG:W07:PDF049:0074
\[
(a+a_1i)^{19}=1,
\]
% ENSEG:W07:PDF049:0075
which, on expanding by Newton's formula and reducing the powers of
$a$, $a_1$, and $i$ by the formulas
% ENSEG:W07:PDF049:0076
\[
\GalSourceError{a^{m(p-1)}=1,\qquad a_1^{m(p-1)}=1}{W07-SE002: these Fermat reductions require a and a_1 to be nonzero modulo p; that restriction is not printed},
\qquad i^3=2,
\]
% ENSEG:W07:PDF049:0077
reduces to
% ENSEG:W07:PDF049:0078
\[
3\bigl[a-a^4a_1^3+(a^5a_1^2+a^2a_1^5)i^2\bigr]=1,
\]
% ENSEG:W07:PDF049:0079
whence, on separating,
% ENSEG:W07:PDF049:0080
\[
3a-3a^4a_1^3=1,\qquad a^5a_1^2+a^2a_1^5=0.
\]

% ENSEG:W07:PDF049:0081
\GalCriticalNote{W07-SE002}{The Fermat reductions require nonzero inputs}{The
printed reductions require $a\ne0$ and $a_1\ne0$ modulo $p$; they fail when
exactly one is zero. The pair chosen immediately afterward,
$a=-1$, $a_1=1$, satisfies those missing conditions, so that computation
remains valid. Any later reuse of the formulas must state the nonzero
hypotheses.}

% ENSEG:W07:PDF049:0082
These last two equations are satisfied by setting $a=-1$, $a_1=1$.
Thus
% ENSEG:W07:PDF049:0083
\[
-1+i
\]
% ENSEG:W07:PDF049:0084
is a primitive root of $x^{19}=1$. Above we found the values $-1$ and
$i$ as primitive roots of $x^2=1$ and $x^{3^2}=1$; it remains only to
multiply together the three quantities
% ENSEG:W07:PDF049:0085
\[
-1,\qquad i,\qquad -1+i,
\]
% ENSEG:W07:PDF049:0086
and the product $i-i^2$ will be a primitive root of the congruence
% ENSEG:W07:PDF049:0087
\[
x^{7^3-1}=1.
\]

% ENSEG:W07:PDF049:0088
Thus here the expression $i-i^2$ has the property that, by raising it
to all powers, one obtains $7^3-1$ distinct expressions of the form
% ENSEG:W07:PDF049:0089
\[
a+a_1i+a_2i^2.
\]

% ENSEG:W07:PDF049:0090
If we wish to obtain the congruence of least degree on which our
primitive root depends, we must eliminate $i$ between the two equations
% ENSEG:W07:PDF049:0091
\[
i^3=2,\qquad \alpha=i-i^2.
\]
% ENSEG:W07:PDF049:0092
We thereby obtain
% ENSEG:W07:PDF049:0093
\[
\alpha^3-\alpha+2=0.
\]

% ENSEG:W07:PDF049:0094
It will be convenient to take as a basis for the imaginaries, and to
represent by $i$, a root of this equation, so that
% ENSEG:W07:PDF049:0095
\[
\tag{$i$}
i^3-i+2=0,
\]

\GalSourcePageBoundary{050}{L0049}{21}{21}

% ENSEG:W07:PDF050:0096
and all the imaginaries of the form
% ENSEG:W07:PDF050:0097
\[
a+a_1i+a_2i^2,
\]
% ENSEG:W07:PDF050:0098
will be obtained by raising $i$ to all powers and reducing by equation
$(i)$.

% ENSEG:W07:PDF050:0099
The principal advantage of the new theory just presented is that it
restores to congruences the property (so useful for ordinary equations)
of admitting exactly as many roots as their degree.

% ENSEG:W07:PDF050:0100
The method for obtaining all these roots will be very simple. First,
\GalSourceError{one can always prepare the given congruence $Fx=0$ so
that it no longer has equal roots, or, in other words, so that it no
longer has a common factor with $F'x=0$, and the means of doing this is
evidently the same as for ordinary equations}{W07-SE003: in characteristic p the derivative can vanish identically, so derivative-gcd reduction alone does not always produce the squarefree part}.

% ENSEG:W07:PDF050:0101
\GalCriticalNote{W07-SE003}{The squarefree method misses inseparable cases}{In
characteristic $p$, $F'$ may vanish identically. For example,
$F=x^3-1=(x-1)^3$ over $\mathbf F_3$ has $F'=0$, and division by
$\gcd(F,F')$ loses the root. A correct squarefree procedure must first include
a $p$th-root extraction branch and then recurse.}

% ENSEG:W07:PDF050:0102
Next, to obtain the integral solutions, it will be enough, as M.~Libri
appears to have been the first to observe, to seek the greatest common
factor of $Fx=0$ and
\GalSourceError{$x^{p-1}=1$}{W07-SE004: this polynomial detects only nonzero residues and omits a possible root x=0}.

% ENSEG:W07:PDF050:0103
\GalCriticalNote{W07-SE004}{The test omits the root zero}{The polynomial
$x^{p-1}-1$ detects only nonzero residue classes. For $F=x$, zero is an
integral solution although the stated greatest common divisor is $1$. To find
all residue roots, one must use $x^p-x$ or test zero separately.}

% ENSEG:W07:PDF050:0104
If one now wishes to obtain the imaginary solutions of the second
degree, one seeks the greatest common factor of $Fx=0$ and
\GalSourceError{$x^{p^2-1}=1$}{W07-SE005: this gcd contains proper-subfield roots and therefore does not isolate roots of exact degree two}; in general, the solutions of order $\nu$ will be given by the
greatest common divisor of $Fx=0$ and
\GalSourceError{$x^{p^\nu-1}=1$}{W07-SE005: this gcd contains every nonzero proper-subfield element whose degree divides nu, rather than isolating exact degree nu}.

% ENSEG:W07:PDF050:0105
\GalCriticalNote{W07-SE005}{The polynomial does not isolate exact degree}{The
stated polynomial contains every nonzero element of each proper subfield whose
degree divides $\nu$. It therefore does not isolate roots of exact degree
$\nu$; when $p=3$ and $\nu=2$, it already contains the degree-one element
$1$. Exact-degree isolation must remove the proper-subfield factors and handle
zero separately.}

% ENSEG:W07:PDF050:0106
It is above all in the theory of permutations, where one constantly
needs to vary the form of the indices, that consideration of the
imaginary roots of congruences appears indispensable. It provides a
simple and easy means of recognizing when a primitive equation is
solvable by radicals, as I shall try to indicate briefly.

% ENSEG:W07:PDF050:0107
Let $fx=0$ be an algebraic equation of degree $p^\nu$; suppose that the
$p^\nu$ roots are denoted by $x_k$, with the index $k$ taking the
$p^\nu$ values determined by the congruence $k^{p^\nu}=k\pmod p$.

% ENSEG:W07:PDF050:0108
Take any rational function $V$ of the $p^\nu$ roots $x_k$. Transform this
function by replacing the index $k$ everywhere with the index
$(ak+b)^{p^r}$, where $a$, $b$, and $r$ are arbitrary constants satisfying
$a^{p^\nu-1}=1$, $b^{p^\nu}=b\pmod p$, and $r$ is an integer.

% ENSEG:W07:PDF050:0109
By giving the constants $a$, $b$, and $r$ all their possible values, one
obtains in all $p^\nu(p^\nu-1)\nu$ ways of

\GalSourcePageBoundary{051}{L0050}{22}{22}

% ENSEG:W07:PDF051:0110
permuting the roots among themselves by substitutions of the form
$[x_k,x_{(ak+b)^{p^r}}]$, and under these substitutions the function $V$
will in general take $p^\nu(p^\nu-1)\nu$ different forms.

% ENSEG:W07:PDF051:0111
Now suppose that the proposed equation $fx=0$ is such that every
function of the roots invariant under the $p^\nu(p^\nu-1)\nu$
permutations just constructed thereby has a rational numerical value.

% ENSEG:W07:PDF051:0112
It follows that, under these circumstances, the equation $fx=0$ will be
solvable by radicals; to reach this conclusion, it is enough to observe
that the value substituted for $k$ in each index can be written in the
three forms
% ENSEG:W07:PDF051:0113
\[
(ak+b)^{p^r}
=\bigl[a(k+b^{1})\bigr]^{p^r}
=a'k^{p^r}+b''
=a'(k+b')^{p^r}.
\]

% ENSEG:W07:PDF051:0114
Those accustomed to the theory of equations will see this without
difficulty.

% ENSEG:W07:PDF051:0115
This observation would be of little importance had I not succeeded in
proving that, conversely, a primitive equation cannot be solvable by
radicals unless it satisfies the conditions just stated. (I except
equations of the ninth and twenty-fifth degrees.)

% ENSEG:W07:PDF051:0125
\GalCriticalNote{POST-P13-A005}{The degree-$9$ and degree-$25$ list is not exhaustive}{Let
$V=\mathbf F_2^4$ and let the irreducible solvable complement
$(C_3\times C_3)\rtimes C_2$ of order $18$ act on $V$. The resulting affine
primitive group has degree $16$ and order $288$, but it is not contained in
$\Gamma L(1,16)$ because $18\nmid60$. The correct general statement is that a
finite solvable primitive group is affine, with a regular elementary-abelian
socle and a faithful irreducible stabilizer. One-dimensional semilinearity
requires an extra hypothesis.}

% ENSEG:W07:PDF051:0116
Thus, for each number of the form $p^\nu$, one can form a group of
permutations such that every function of the roots invariant under these
permutations must have a rational value whenever the equation of degree
$p^\nu$ is primitive and solvable by radicals.

% ENSEG:W07:PDF051:0117
Moreover, only equations of such a degree $p^\nu$ are both primitive
and solvable by radicals.

% ENSEG:W07:PDF051:0118
The general theorem just stated makes precise and develops the
conditions I gave in the April issue of the \emph{Bulletin}. It gives a
means of forming a function of the roots whose value will be rational
whenever the primitive equation of degree $p^\nu$ is solvable by
radicals, and consequently leads to criteria for the solvability of
these equations by calculations which, if not practicable, are at least
possible in theory.

% ENSEG:W07:PDF051:0119
It should be noted that, when $\nu=1$, the various values of $k$ are
nothing other than the sequence of integers. The

\GalSourcePageBoundary{052}{L0051}{23}{23}

% ENSEG:W07:PDF052:0120
substitutions of the form $(x_k,x_{ak+b})$ will number $p(p-1)$.

% ENSEG:W07:PDF052:0121
The function that, in the case of equations solvable by radicals, must
have a rational value will in general depend on an equation of degree
$1\mathbin{\cdot}2\mathbin{\cdot}3\cdots(p-2)$, to which the rational-root
method must consequently be applied.

\GalArticleEndOrnament{figures/PDF052_L0051_TERMINAL_ORNAMENT_ENHANCED.png}
% ENSEG:W07:PDF052:0123
\GalSourceImage{052}{L0051}{terminal printer ornament; raw and conservatively enhanced crops retained}{figures/PDF052_L0051_TERMINAL_ORNAMENT_RAW.png}{figures/PDF052_L0051_TERMINAL_ORNAMENT_ENHANCED.png}

\GalSourceBlankPage{053}{L0052}{24}{not visibly printed}
